\RequirePackage{plautopatch}
\documentclass[a4paper]{ltjsarticle}
\usepackage{tcolorbox,amsmath,amssymb,graphicx,ascmac,fancybox,framed,tikz,comment,calc}
\tcbuselibrary{breakable, skins, theorems}
\usetikzlibrary{positioning, intersections, calc, arrows.meta,math}
\usepackage[top=20truemm,left=15truemm,right=15truemm,bottom=16truemm]{geometry}
%
% \white :計算用余白の挟み込み
\newcommand{\white}{\newpage\leavevmode
\begin{center}
    \noindent(計算用余白)
\end{center}}
% \exc[サイズ]{数値} :a\exc{b}でa^b, 指数に指数を載せる時に使用　数式環境内ではサイズ0.95
\newcommand{\exc}[2][0.75]{^{\textstyle\scalebox{#1}[#1]{\raise0.25ex\hbox{$#2$}}}}
% \exfrac[サイズ]{分母}{分子} :指数部分に分数を持ってくるときに使う．積分範囲などに使う場合サイズ0.75~0.8を推奨
\newcommand{\exfrac}[3][0.90]{\tfrac{\lower0.1ex\hbox{\scalebox{#1}[#1]{$#2$}}}{\raise0.3ex\scalebox{#1}[#1]{$#3$\rule{0pt}{7.5pt}}}}
% \dint[文字サイズ]{下付}{上付} :インテグラルのdisplay省略、積分範囲の表示サイズを変更、
\newcommand{\dint}[3][0.75]{\displaystyle\int_{\scalebox{#1}[#1]{$#2$}}^{\scalebox{#1}[#1]{$#3$}}}
% \dsqrt[n乗根]{中身}　:√を重ねる時用
\newcommand{\dsqrt}[2][]{\displaystyle\sqrt[\raise0ex\hbox{{\scalebox{0.625}[0.625]{$#1$}}}]{#2}}
% \dlim{下に付けるヤツ} :文字サイズ変更に適用させたlim
\newcommand{\dlim}[1]{\raise0.3ex\hbox{\scalebox{1.04}{${\normalsize \displaystyle\lim_{#1}}$}}\,}
% \cmb{n}{r}  \prm{n}{r}　: nCr nPr
\def\cmb#1#2{_{#1}\mathrm{C}_{#2}}\def\prm#1#2{_{#1}\mathrm{P}_{#2}}
% \w :(n)の後の空白用
\def\w{\;\;\;}\def\spw{\;\;\,\,}
% \hs \vs \vhs  :未使用
\def\hs{\hspace{1.38cm}}\def\vs{\vspace{-2mm}\\[1mm]}\def\vhs{\vs\hs}
\def\mrm{\mathrm}
\def\ve{\overrightarrow}
\def\linera{\\[-0.4mm]｜\\[-0.5mm]｜}
\def\linerb{\linera\linera}
\def\linerc{\linerb\linerb}
\def\linerd{\linerc\linerc}
\def\linere{\linerd\linerd\linerc\linerb}

% 表紙のページ番号を省略
\pagestyle{empty}

\begin{document}
\fontsize{14truept}{14truept}\selectfont\noindent


% p1
\noindent
模\;擬\;試\;験\spw 数\;学\spw 解\;答\;用\;紙 \spw(\;1\;枚\,目\;表\;)\\[1mm]
{\small 解答は枠内に収めること．\scalebox{1.1}{\fbox{2}}\,の解答欄は裏面にある．}
\begin{framed}\noindent
    \scalebox{1.5}{\fbox{1}}\begin{center}
        ｜ \linere\
    \end{center}
\end{framed}\leavevmode\\[-10mm]
\,\underbar{　\hspace{17cm}　}\\[2mm]
\noindent
$\raise2.75ex\hbox{\scalebox{0.8}{
    \begin{tabular}{|l|r|} \hline
    \;\;{\small フリガナ} & \hspace{5cm} \\ \hline
    \lower1ex\hbox{　\hspace{-0.1cm}}{\large 氏\; 名}\raise1ex\hbox{　\hspace{-0.1cm}}& \hspace{9cm} \\ \hline 
        \end{tabular}}}$
\hfill $\raise1.2ex\hbox{{\HUGE \fbox{　\hspace{-0.9cm} {\raise0.4ex\hbox{\large 採\,点}}　\hspace{-0.595cm}}\fbox{　{\small /}　}}}$


% p2
\newpage\noindent
$模\;擬\;試\;験\spw 数\;学\spw 解\;答\;用\;紙 \spw(\;1\;枚\,目\;裏\;)$ \\[2mm]
{\small 解答は枠内に収めること．}
\begin{framed}\noindent
    \scalebox{1.5}{\fbox{2}}\begin{center}
        ｜ \linere\
    \end{center}
\end{framed}\leavevmode\\[-10mm]
\,\underbar{　\hspace{17cm}　}\\[2mm]
\noindent
\hfill $\raise1.2ex\hbox{{\HUGE \fbox{　\hspace{-0.9cm} {\raise0.4ex\hbox{\large 採\,点}}　\hspace{-0.595cm}}\fbox{　{\small /}　}}}$

% p3
\newpage\noindent
$模\;擬\;試\;験\spw 数\;学\spw 解\;答\;用\;紙 \spw(\;2\;枚\,目\;表\;)$\\[1mm]
{\small 解答は枠内に収めること．\scalebox{1.1}{\fbox{4}}\,の解答欄は裏面にある．}
\begin{framed}\noindent
    \scalebox{1.5}{\fbox{3}}\begin{center}
        ｜ \linere\
    \end{center}
\end{framed}\leavevmode\\[-10mm]
\,\underbar{　\hspace{17cm}　}\\[2mm]
\noindent
$\raise2.75ex\hbox{\scalebox{0.8}{
    \begin{tabular}{|l|r|} \hline
    \;\;{\small フリガナ} & \hspace{5cm} \\ \hline
    \lower1ex\hbox{　\hspace{-0.1cm}}{\large 氏\; 名}\raise1ex\hbox{　\hspace{-0.1cm}}& \hspace{9cm} \\ \hline 
        \end{tabular}}}$
\hfill $\raise1.2ex\hbox{{\HUGE \fbox{　\hspace{-0.9cm} {\raise0.4ex\hbox{\large 採\,点}}　\hspace{-0.595cm}}\fbox{　{\small /}　}}}$





% p4
\newpage\noindent
$模\;擬\;試\;験\spw 数\;学\spw 解\;答\;用\;紙 \spw(\;2\;枚\,目\;裏\;)$ \\[2mm]
{\small 解答は枠内に収めること．}
\begin{framed}\noindent
    \scalebox{1.5}{\fbox{4}}\begin{center}
        ｜ \linere\
    \end{center}
\end{framed}\leavevmode\\[-10mm]
\,\underbar{　\hspace{17cm}　}\\[2mm]
\noindent
\hfill $\raise1.2ex\hbox{{\HUGE \fbox{　\hspace{-0.9cm} {\raise0.4ex\hbox{\large 採\,点}}　\hspace{-0.595cm}}\fbox{　{\small /}　}}}$


% p5
\newpage\noindent
$模\;擬\;試\;験\spw 数\;学\spw 解\;答\;用\;紙 \spw(\;3\;枚\,目\;表\;)　$\\[1mm]
{\small 解答は枠内に収めること．\scalebox{1.1}{\fbox{6}}\,の解答欄は裏面にある．}
\begin{framed}\noindent
    \scalebox{1.5}{\fbox{5}}\begin{center}
        ｜ \linere\
    \end{center}
\end{framed}\leavevmode\\[-10mm]
\,\underbar{　\hspace{17cm}　}\\[2mm]
\noindent
$\raise2.75ex\hbox{\scalebox{0.8}{
    \begin{tabular}{|l|r|} \hline
    \;\;{\small フリガナ} & \hspace{5cm} \\ \hline
    \lower1ex\hbox{　\hspace{-0.1cm}}{\large 氏\; 名}\raise1ex\hbox{　\hspace{-0.1cm}}& \hspace{9cm} \\ \hline 
        \end{tabular}}}$
\hfill $\raise1.2ex\hbox{{\HUGE \fbox{　\hspace{-0.9cm} {\raise0.4ex\hbox{\large 採\,点}}　\hspace{-0.595cm}}\fbox{　{\small /}　}}}$

%

% p6
\newpage\noindent
$模\;擬\;試\;験\spw 数\;学\spw 解\;答\;用\;紙 \spw(\;3\;枚\,目\;裏\;)$ \\[2mm]
{\small 解答は枠内に収めること．}
\begin{framed}\noindent
    \scalebox{1.5}{\fbox{6}}\begin{center}
        ｜ \linere\
    \end{center}
\end{framed}\leavevmode\\[-10mm]
\,\underbar{　\hspace{17cm}　}\\[2mm]
\noindent
\hfill $\raise1.2ex\hbox{{\HUGE \fbox{　\hspace{-0.9cm} {\raise0.4ex\hbox{\large 採\,点}}　\hspace{-0.595cm}}\fbox{　{\small /}　}}}$
%



%
%
%
\end{document}
